\documentclass[10pt,a4paper]{amsart}
\usepackage[margin=19mm]{geometry}
\usepackage{amsmath,amssymb,mathtools}
\usepackage[scr=rsfs,cal=euler]{mathalfa}
\usepackage{xfrac}
\usepackage[unicode,hidelinks,bookmarks=false]{hyperref}
\newcommand{\ie}{i.e.\ }
\newcommand{\cf}{cf.\ }
\newcommand{\resp}{respectively\ }
\newcommand{\Subsection}{\subsection}
\providecommand{\nobreakdash}{\nobreak\hbox{-}}
\setlength{\emergencystretch}{2em}
\multlinegap0pt
\usepackage{enumerate}
\usepackage{amssymb}
\usepackage{tikz}
\usepackage{mathtools}
\newtheorem{theorem}{Theorem}
\title{This is the title}
\author{K. Mahesh Krishna}
\date{Source: arXiv:2304.03324v2 (CC BY 4.0)}
\begin{document}
\maketitle

\noindent\emph{Source and licence.} This is the title. Version arXiv:2304.03324v2 (CC BY 4.0), distributed by arXiv under the Creative Commons Attribution 4.0 International licence, http://creativecommons.org/licenses/by/4.0/, as recorded in the arXiv licence metadata for this exact version and shown on the abstract page https://arxiv.org/abs/2304.03324v2. Full licence notice: \texttt{licenses/arxiv-2304.03324-CC-BY-4.0.txt}.\par\smallskip

\noindent\textit{Editorial scope.} Counted unit: the article's theorem FDS (source0.tex lines 184--189). The article's complete original proof of it is delivered with it (source0.tex lines 190--224, one \texttt{proof} environment of 35 lines). No mathematics is written, repaired or re-derived: the statement and the proof are the article's own lines, in the article's own notation, and no retained line is substituted or re-flowed.  No further source-local context is needed: every cross-reference of every retained line resolves inside the two delivered spans.  Nothing was omitted from any retained span, so the package carries no omission record.  No retained line contains a citation, so the wrapper carries no bibliography and no bibliography entry.  No retained line contains a figure environment, an image-inclusion command or drawing code, so no image is delivered and the package declares no asset dependency.  Editorial formatting only, in addition: an AMS \texttt{amsart} preamble that restates the article's own declarations and macros used by the retained lines, namely the environments \texttt{theorem} on separate counters, together with the standard packages those lines need.  The retained environments print in this document as Theorem 1 in order of appearance, whereas the article prints the same items with the numbering of its own full document.  The complete native source of the article as distributed in the arXiv e-print for this exact version is bundled unchanged as \texttt{sources/139792/source0.tex}.
\begin{theorem} \label{FDS}(\textbf{Functional Donoho-Stark-Elad-Bruckstein-Ricaud-Torr\'{e}sani Uncertainty Principle}) 
Let $(\{f_j\}_{j=1}^n, \{\tau_j\}_{j=1}^n)$ and $(\{g_k\}_{k=1}^m, \{\omega_k\}_{k=1}^m)$ be  p-Schauder frames  for a Banach space $\mathcal{X}$. Then for every $x \in \mathcal{X}\setminus\{0\}$,  we have 
	\begin{align}\label{FDSUNCE}
		\|\theta_f x\|_0^\frac{1}{p}\|\theta_g x\|_0^\frac{1}{q} \geq 	\frac{1}{\displaystyle\max_{1\leq j\leq n, 1\leq k\leq m}|f_j(\omega_k)|}\quad \text{and} \quad \|\theta_g x\|_0^\frac{1}{p}\|\theta_f x\|_0^\frac{1}{q}\geq \frac{1}{\displaystyle\max_{1\leq j\leq n, 1\leq k\leq m}|g_k(\tau_j)|}.
	\end{align}
\end{theorem}

\begin{proof}
	Let $x \in \mathcal{X}\setminus\{0\}$ and $q$ be the conjugate index of $p$. First using $\theta_f$ is an isometry and later using $\theta_g$ is an isometry, we get 
	\begin{align*}
\|x\|^p&=\|\theta_fx\|^p=\sum_{j=1}^n|f_j(x)|^p=\sum_{j \in \operatorname{supp}(\theta_fx)}|f_j(x)|^p\\
&=\sum_{j\in \operatorname{supp}(\theta_fx)}\left|f_j\left(\sum_{k=1}^mg_k(x)\omega_k\right)\right|^p
=\sum_{j \in \operatorname{supp}(\theta_fx)}\left|\sum_{k=1}^mg_k(x)f_j(\omega_k)\right|^p\\
&=\sum_{j \in \operatorname{supp}(\theta_fx)}\left|\sum_{k \in  \operatorname{supp}(\theta_gx)}g_k(x)f_j(\omega_k)\right|^p\leq \sum_{j \in \operatorname{supp}(\theta_fx)}\left(\sum_{k \in  \operatorname{supp}(\theta_gx)}|g_k(x)f_j(\omega_k)|\right)^p\\
&\leq \left(\displaystyle\max_{1\leq j\leq n, 1\leq k\leq m}|f_j(\omega_k)|\right)^p\sum_{j \in \operatorname{supp}(\theta_fx)}\left(\sum_{k \in  \operatorname{supp}(\theta_gx)}|g_k(x)|\right)^p\\
&=\left(\displaystyle\max_{1\leq j\leq n, 1\leq k\leq m}|f_j(\omega_k)|\right)^p \|\theta_f x\|_0\left(\sum_{k \in  \operatorname{supp}(\theta_gx)}|g_k(x)|\right)^p\\
&\leq \left(\displaystyle\max_{1\leq j\leq n, 1\leq k\leq m}|f_j(\omega_k)|\right)^p \|\theta_f x\|_0\left(\sum_{k \in  \operatorname{supp}(\theta_gx)}|g_k(x)|^p\right)^\frac{p}{p}\left(\sum_{k \in  \operatorname{supp}(\theta_gx)}1^q\right)^\frac{p}{q}\\
&=\left(\displaystyle\max_{1\leq j\leq n, 1\leq k\leq m}|f_j(\omega_k)|\right)^p \|\theta_f x\|_0\|\theta_g x\|^p\|\theta_g x\|_0^\frac{p}{q}\\
&=\left(\displaystyle\max_{1\leq j\leq n, 1\leq k\leq m}|f_j(\omega_k)|\right)^p \|\theta_f x\|_0\|x\|^p\|\theta_g x\|_0^\frac{p}{q}.
	\end{align*}
Therefore 
\begin{align*}
	\frac{1}{\displaystyle\max_{1\leq j\leq n, 1\leq k\leq m}|f_j(\omega_k)|}\leq \|\theta_f x\|_0^\frac{1}{p}\|\theta_g x\|_0^\frac{1}{q}.
\end{align*}
On the other way, first using $\theta_g$ is an isometry and $\theta_f$ is an isometry, we get

\begin{align*}
\|x\|^p&=\|\theta_gx\|^p=\sum_{k=1}^m|g_k(x)|^p=\sum_{k \in \operatorname{supp}(\theta_gx)}|g_k(x)|^p\\
&=\sum_{k\in \operatorname{supp}(\theta_gx)}\left|g_k\left(\sum_{j=1}^nf_j(x)\tau_j\right)\right|^p
=\sum_{k \in \operatorname{supp}(\theta_gx)}\left|\sum_{j=1}^nf_j(x)g_k(\tau_j)\right|^p\\
&=\sum_{k \in \operatorname{supp}(\theta_gx)}\left|\sum_{j\in \operatorname{supp}(\theta_fx)}f_j(x)g_k(\tau_j)\right|^p\leq \sum_{k \in \operatorname{supp}(\theta_gx)}\left(\sum_{j\in \operatorname{supp}(\theta_fx)}|f_j(x)g_k(\tau_j)|\right)^p\\
&\leq \left(\displaystyle\max_{1\leq j\leq n, 1\leq k\leq m}|g_k(\tau_j)|\right)^p\sum_{k \in \operatorname{supp}(\theta_gx)}\left(\sum_{j\in \operatorname{supp}(\theta_fx)}|f_j(x)|\right)^p\\
&=\left(\displaystyle\max_{1\leq j\leq n, 1\leq k\leq m}|g_k(\tau_j)|\right)^p\|\theta_g x\|_0\left(\sum_{j\in \operatorname{supp}(\theta_fx)}|f_j(x)|\right)^p\\
&\leq \left(\displaystyle\max_{1\leq j\leq n, 1\leq k\leq m}|g_k(\tau_j)|\right)^p\|\theta_g x\|_0\left(\sum_{j\in \operatorname{supp}(\theta_fx)}|f_j(x)|^p\right)^\frac{p}{p}\left(\sum_{j\in \operatorname{supp}(\theta_fx)}1^q\right)^\frac{p}{q}\\
&=\left(\displaystyle\max_{1\leq j\leq n, 1\leq k\leq m}|g_k(\tau_j)|\right)^p\|\theta_g x\|_0\|\theta_f x\|^p\|\theta_f x\|_0^\frac{p}{q}\\
&=\left(\displaystyle\max_{1\leq j\leq n, 1\leq k\leq m}|g_k(\tau_j)|\right)^p\|\theta_g x\|_0\|x\|^p\|\theta_f x\|_0^\frac{p}{q}.
\end{align*}
Therefore 
\begin{align*}
\frac{1}{\displaystyle\max_{1\leq j\leq n, 1\leq k\leq m}|g_k(\tau_j)|}\leq \|\theta_g x\|_0^\frac{1}{p}\|\theta_f x\|_0^\frac{1}{q}.
\end{align*}
\end{proof}

\end{document}
